\documentclass[12pt,a4paper]{article}
\usepackage[UTF8]{ctex}
\usepackage{geometry}
\usepackage{setspace}
\usepackage{titlesec}
\usepackage{fancyhdr}
\usepackage{amsmath}
\usepackage{graphicx}
\usepackage{hyperref}
\usepackage{booktabs}
\usepackage{enumitem}

\geometry{left=2.5cm,right=2.5cm,top=2.5cm,bottom=2.5cm}
\onehalfspacing

\pagestyle{fancy}
\fancyhf{}
\fancyhead[L]{马歇尔计划与亲密关系}
\fancyhead[R]{\thepage}
\renewcommand{\headrulewidth}{0.4pt}

\titleformat{\section}{\Large\bfseries}{\thesection}{1em}{}
\titleformat{\subsection}{\large\bfseries}{\thesubsection}{1em}{}

\title{\textbf{骨头与鞭子：从马歇尔计划到亲密关系的权力美学}}
\author{匿名作者 \\ \small{量化龙虾研究院}}
\date{2026年3月}

\begin{document}

\maketitle

\begin{abstract}
本文通过深度剖析二战后美国对西德与日本的政治经济改造工程，提炼出一套"恩威并施"的权力运作逻辑。研究认为，这一历史范式与亲密关系中的权力动力学存在惊人的同构性。通过考察马歇尔计划中的经济援助（骨头）与政治控制（鞭子）的双重机制，本文建立了一个跨学科的分析框架，将地缘政治智慧转化为恋爱哲学。研究发现，成功的长期关系管理需要同时满足对方的安全感需求与成长性需求，这种平衡艺术与战后重建的政治智慧异曲同工。

\textbf{关键词：} 马歇尔计划；亲密关系；权力动力学；依附理论；博弈论
\end{abstract}

\newpage
\tableofcontents
\newpage

\section{引言：一个危险的类比}

\subsection{研究缘起}

1947年6月5日，哈佛大学的毕业典礼上，美国国务卿乔治·马歇尔发表了一场改变世界格局的演讲。他宣布："我们的政策不是针对任何国家或主义，而是针对饥饿、贫穷、绝望和混乱。"这场后来被称为"马歇尔计划"的演说，拉开了人类历史上最大规模的国际援助工程的序幕。

然而，当我们剥开人道主义的修辞外衣，深入检视这一工程的运作机制时，会发现一个更为复杂的权力图景：西德和日本在获得空前经济援助的同时，也接受了前所未有的政治控制。经济复兴与主权让渡同步进行，物质馈赠与制度改造并行不悖。这种看似矛盾的操作，实际上构成了一套精妙的权力技术——我们不妨称之为"骨头与鞭子"的治理术。

本文试图论证一个看似离经的命题：这套诞生于冷战地缘政治的权力技术，对于理解当代亲密关系中的权力动力学具有出人意料的启发价值。换言之，美国如何"驯服"战败国，与一个人如何经营长期亲密关系，在深层逻辑上存在结构性的相似。

\subsection{问题的理论意义}

这一研究具有三重理论意义：

首先，它打破了国际政治与微观社会学之间的学科壁垒。传统上，国际关系学者关注的是国家间的权力博弈，而亲密关系研究者聚焦于个体间的情感互动。本文认为，这两种看似悬殊的实践领域，实际上共享着相似的权力语法。

其次，它为亲密关系研究引入了历史制度的分析视角。现有的恋爱心理学多基于问卷调查和实验研究，缺乏对权力运作的长时段历史考察。通过回溯马歇尔计划这一经典案例，我们可以更清晰地识别出可持续关系中权力平衡的必要条件。

最后，它为理解"健康的控制"提供了一个概念框架。在后现代语境下，"控制"已成为一个负面词汇。但本文认为，某些形式的控制——如果是以促进对方成长为导向的——可以成为关系的建设性要素。

\section{马歇尔计划：权力运作的双重逻辑}

\subsection{历史背景：战后废墟中的机遇}

1945年，欧洲满目疮痍。战争摧毁了大部分工业基础设施，饥荒和疾病在蔓延。在西德，工业产量仅为1938年的27\%；在日本，城市面积的40\%被夷为平地，工业产量下降到战前的10\%以下。这种毁灭性的物质条件，既构成了人道主义危机，也提供了政治重构的历史机遇。

美国面临着一个根本性的战略问题：如何防止这两个战败国再次成为战争策源地，同时又不让它们落入苏联的势力范围？传统的处置方式——如一战后的赔款和羞辱——已被证明适得其反。美国决策者意识到，必须发明一种新的权力技术，既能控制潜在的威胁，又能赢得被统治者的认同。

\subsection{"骨头"：经济援助的给予逻辑}

马歇尔计划的核心是经济援助。从1948年到1952年，美国向西欧提供了约130亿美元（相当于今天的1500亿美元）的援助，其中西德获得了约14亿美元。日本虽然不在马歇尔计划范围内，但也通过其他渠道获得了约22亿美元的援助。

这种援助不是简单的馈赠，而是一种精心设计的"给予"：

\textbf{1. 条件性给予}

援助不是无条件的。接受国必须同意货币稳定、平衡预算、拆除贸易壁垒等一系列经济政策改革。换言之，经济援助本身就是一种权力工具——你要我的钱，就得听我的话。这种"条件性"确保了援助不仅是物质馈赠，更是制度改造的杠杆。

\textbf{2. 选择性给予}

援助是有选择的。工业设备、粮食、原材料优先供给，而军事相关技术和战略资源则严格管制。这种选择性确保了受援国的经济发展方向符合美国的战略利益——你可以富，但不能强到威胁我。

\textbf{3. 象征性给予}

援助不仅是物质的，也是象征的。马歇尔计划被包装成"自由世界"对"共产主义威胁"的回应，援助本身就是一种意识形态表态。接受援助，就等于认同了美国的价值观和战略目标。

\subsection{"鞭子"：政治控制的约束逻辑}

与经济援助并行的是政治控制。这种控制采取了多种形式：

\textbf{1. 军事存在}

美国在西德和日本建立了大规模的军事基地。这些基地不仅是对苏联的威慑，也是对东道国的控制。只要美国的坦克还在柏林郊外、战斗机还在冲绳起降，这两个国家就不可能脱离美国的轨道。军事存在是一种无声但永不休止的"鞭子"。

\textbf{2. 制度设计}

美国深度参与了西德和日本的宪法制定。西德的基本法确立了联邦制和议会民主制，明确限制了中央政府的权力；日本的和平宪法更是将"放弃战争"写入条款，从法律层面解除了日本的武装能力。这种制度设计确保了两国即使恢复主权，也无法再次成为军事威胁。

\textbf{3. 精英培养}

美国通过奖学金、培训项目、交流计划，培养了整整一代亲美的政治、经济和文化精英。西德的阿登纳、艾哈德，日本的吉田茂、池田勇人，都与美国有着密切的联系。这种"精英俘获"确保了即使政治人物更替，政策的连续性也能得到保证。

\section{案例比较：西德与日本的驯化艺术}

\subsection{西德案例：欧洲的"模范生"}

\subsubsection{骨头：从废墟到经济奇迹}

战后西德的经济复兴被称为"经济奇迹"（Wirtschaftswunder）。1948年的货币改革、马歇尔计划的援助、以及艾哈德推行的"社会市场经济"政策，共同创造了这一奇迹。从1948年到1960年，西德的GDP年均增长率达到8\%，失业率从11\%下降到1\%，工业产量超过了战前水平。

这种经济成功有几个关键因素：马歇尔计划提供了急需的资本货物和原材料；美国向西德产品开放市场；美国企业对西德进行了大规模的技术授权和直接投资。这些"骨头"让西德重新站了起来。

\subsubsection{鞭子：看不见的缰绳}

然而，与经济援助相伴的是严格的政治控制：西德被完全解除武装，没有军队，没有军工产业；即使在恢复主权后，西德的军事力量也被纳入北约框架，实际上受美国指挥；美国在柏林保留了军事存在，这是对西德的持续提醒——你的安全依赖于我的保护。

\subsection{日本案例：东亚的"优等生"}

\subsubsection{骨头：从焦土到经济大国}

日本的重建比西德更为彻底。1945年的日本，城市废墟、工业崩溃、粮食短缺。但到了1960年代，日本已经成为世界第三大经济体。美国对日本的"骨头"包括直接援助、朝鲜战争红利、市场准入、技术授权和金融支持。

\subsubsection{鞭子：和平宪法的牢笼}

然而，与慷慨援助相伴的是严格约束：和平宪法第九条明确规定日本放弃战争能力；美国在日本建立了数百个军事基地；日美安保条约确立了日本对美国的单方面安全依赖；GHQ直接干预日本政治改革。

\subsection{比较分析：驯化的共同逻辑}

西德和日本的案例虽然具体条件不同，但共享着相似的驯化逻辑：

\begin{table}[h]
\centering
\begin{tabular}{lcc}
\toprule
\textbf{维度} & \textbf{西德} & \textbf{日本} \\
\midrule
经济援助规模 & 14亿美元 & 22亿美元 \\
军事存在 & 北约框架 & 安保条约 \\
制度约束 & 基本法 & 和平宪法 \\
结果 & 依附性繁荣 & 半主权繁荣 \\
\bottomrule
\end{tabular}
\caption{西德与日本驯化机制比较}
\end{table}

这种共同逻辑的核心是：通过经济援助建构起受援国对美国的结构性依赖；通过军事存在和制度设计将外部约束转化为内在认同；允许经济发展但控制发展方向；通过援助带来的经济成果生产出美国领导的合法性。

\section{理论提炼：权力美学的六个维度}

通过上述案例分析，我们可以提炼出"骨头与鞭子"权力技术的六个维度：

\subsection{维度一：结构性依赖的建构}

权力的基础不是武力，而是依赖。马歇尔计划的精髓在于，它建构了一种结构性的、难以替代的依赖关系。西德需要美国的市场和技术，日本需要美国的安全保护和市场准入。这种依赖不是偶然的、暂时的，而是嵌入到经济和社会结构之中的。

\textbf{理论意义}：持久的权力关系必须建立在结构性依赖之上。短期的影响力可以依靠武力或贿赂，但长期的支配需要让被统治者产生对统治者的深层需要。

\subsection{维度二：不对称互惠的维持}

马歇尔计划不是单向的馈赠，而是一种不对称的互惠。美国提供经济援助，西德和日本提供政治合作和市场开放。双方都从中受益，但美国获得的是战略主导权，而西德和日本获得的是经济发展权。这种不对称性确保了美国的核心利益——战略控制——得到保障。

\textbf{理论意义}：可持续的权力关系必须包含某种形式的互惠。纯粹的剥削会引发反抗，但互惠的形式和程度可以是不对称的。

\subsection{维度三：自主性的边界控制}

美国并不试图控制西德和日本的一切。相反，它给予这两个国家相当大的经济自主权，甚至在某些领域（如产业政策）鼓励它们发展自己的特色。美国控制的是边界——军事、外交、意识形态——而将经济和社会发展的空间留给当地。

\textbf{理论意义}：过度控制会窒息创造性和主动性，最终削弱被统治者的价值。明智的权力运作是控制边界，而非控制一切。

\subsection{维度四：合法性的持续生产}

权力需要合法性。马歇尔计划的合法性来自于它带来的经济成果。当西德和日本的人民享受到前所未有的物质生活改善时，他们会自然地认同这一套政治安排。经济增长不仅是物质馈赠，更是合法性资源。

\textbf{理论意义}：持久的权力关系必须持续生产合法性。最有效的合法性生产机制是让对方体验到合作带来的实际利益。

\subsection{维度五：抵抗能力的系统削弱}

美国系统性地削弱了西德和日本的抵抗能力。解除武装、制度约束、精英俘获、情报渗透——这些措施确保了即使出现不满，也难以转化为有效的政治挑战。更重要的是，这些措施很多是隐蔽的、制度化的，被统治者甚至意识不到自己在被控制。

\textbf{理论意义}：明智的权力运作不是压制抵抗，而是预防抵抗。最完美的控制是让被统治者意识不到自己在被控制。

\subsection{维度六：长期主义的时间视野}

马歇尔计划不是短期操作，而是长期投资。美国从1948年开始提供援助，直到1960年代才看到完整的回报。这种长期主义的时间视野，使得美国能够容忍短期的成本，追求长期的收益。

\textbf{理论意义}：可持续的权力关系需要长期主义。真正的权力艺术在于，让对方相信这种关系是持久的、值得信任的。

\section{从地缘政治到亲密关系：模型的应用}

现在，我们可以将这一分析框架应用于亲密关系。我们需要澄清：这一类比不是要将亲密关系简化为权力斗争，而是要识别出其中相似的权力动力学。健康的亲密关系不是没有权力，而是权力得到了正当的、建设性的运作。

\subsection{关系中的"结构性依赖"}

在亲密关系中，结构性依赖体现为：情感支持、经济保障、社会认同、子女抚养等多个维度。一个人之所以留在一段关系中，是因为对方满足了某些难以替代的需要。

\textbf{对照马歇尔计划}：美国让西德和日本离不开它的市场、技术和安全保护。同样，在亲密关系中，让对方在某些关键维度上依赖你，是关系持久的基础。

\textbf{实践意义}：成为对方生活中不可或缺的一部分；在某些领域建立不可替代的优势；避免让对方感到你随时可以被替代。

\subsection{关系中的"不对称互惠"}

亲密关系是一种互惠关系，但这种互惠可以是不对称的。一方可能提供更多的经济资源，另一方提供更多的情感劳动。关键在于，双方都感到交易是公平的，即使在客观上是不对称的。

\textbf{对照马歇尔计划}：美国提供经济援助，西德和日本提供政治合作。双方都受益，但美国保持了战略主导权。

\textbf{实践意义}：确保双方都能从关系中获得核心利益的满足；不追求每一项交换的绝对对称，但追求长期的大致平衡；让对方感到关系是公平的。

\subsection{关系中的"边界控制"与"自主空间"}

健康的亲密关系给予对方充分的自主空间。你不需要控制对方的每一个决定、每一个想法。你控制的是边界——忠诚、透明、共同目标——而在边界之内，对方有充分的自由。

\textbf{对照马歇尔计划}：美国控制西德和日本的军事和外交，但给予它们经济和社会发展的自主权。

\textbf{实践意义}：设定清晰的关系边界；在边界之内给予对方充分的自由和尊重；不进行微观管理，让对方保持独立人格。

\subsection{关系中的"合法性生产"}

为什么对方愿意和你在一起？不仅因为你满足了某些需要，更因为这段关系让对方体验到了成长和幸福。这种正向体验就是关系的合法性来源。

\textbf{对照马歇尔计划}：马歇尔计划的合法性来自于经济增长带来的物质改善。

\textbf{实践意义}：持续为对方创造正向体验；帮助对方成长为更好的自己；让对方感到和你在一起是值得的。

\subsection{关系中的"建设性张力"}

马歇尔计划中的"鞭子"不是纯粹的压迫，而是"建设性的压力"。美国通过制度约束和军事存在，确保西德和日本不会偏离轨道。但这种约束也创造了稳定性，使这两个国家得以专注于经济发展。

同样，在亲密关系中，某些形式的"压力"可以是建设性的。比如，要求对方保持健康、追求成长、承担责任。这些压力不是压迫，而是对共同未来的投资。

\textbf{实践意义}：对对方有合理的期望和要求；以促进成长为导向施加压力，而非以控制为目的；让"压力"服务于双方的长期利益。

\subsection{关系中的"长期主义"}

真正的关系不是短期的激情，而是长期的承诺。马歇尔计划的成功在于美国展现出长期主义的视野——它愿意在短期内投入资源，追求长期的战略收益。

同样，亲密关系中的权力运作需要长期主义。不是追求短期的控制或利益，而是投资于关系的长期健康。

\textbf{实践意义}：不追求短期利益最大化，而追求关系的长期可持续；展现出对关系的长期承诺；让对方相信这段关系是值得长期投资的。

\section{实践指南：关系中的"马歇尔计划"}

基于上述理论框架，我们可以提出一套亲密关系中的"马歇尔计划"——如何以建设性的方式运用权力，促进双方的成长和关系的持久。

\subsection{给予的艺术：如何投喂"骨头"}

\subsubsection{识别对方的核心需求}

每个人在亲密关系中的核心需求不同。有的人需要情感支持，有的人需要经济保障，有的人需要社会认同，有的人需要智力刺激。给予的第一步是准确识别对方的核心需求。

\textbf{实践建议}：观察对方在什么情况下最感谢你；直接询问对方最需要什么；区分"想要"和"需要"。

\subsubsection{创造不可替代的价值}

单纯的给予不够，关键是创造不可替代的价值。如果对方可以轻易从别处获得同样的东西，你的给予就不会产生结构性依赖。

\textbf{实践建议}：发展只有你能提供的独特价值；建立共同的历史和记忆；持续投资于自己的成长。

\subsubsection{有条件但不是交易}

马歇尔计划的援助是有条件的，但这些条件服务于长期目标，而非短期交易。同样，关系中的给予可以有期望，但不应沦为交易。

\textbf{实践建议}：让对方知道你的期望，但不要每次给予都立即要求回报；关注长期互惠，而非短期交换。

\subsection{约束的艺术：如何挥舞"鞭子"}

\subsubsection{设定健康的边界}

"鞭子"的第一层含义是边界。边界不是控制，而是关系的护栏。比如，什么行为是不可接受的，什么是必须透明的，什么是共同的底线。

\textbf{实践建议}：明确沟通你的边界；边界应当是合理的、可执行的；一致性是关键。

\subsubsection{施加建设性的压力}

"鞭子"的第二层含义是成长压力。在某些情况下，你需要推动对方走出舒适区，承担更多责任，追求更高目标。这种压力不是压迫，而是对对方潜力的信任。

\textbf{实践建议}：以促进成长为目的，而非以控制为目的；压力应当是适度的；在施加压力的同时提供支持。

\subsubsection{维持适当的不对称}

"鞭子"的第三层含义是保持一定的优势。这不意味着傲慢或压迫，而是确保你在关系中有足够的筹码，不至于完全被动。

\textbf{实践建议}：保持经济、社交或情感上的相对独立；不要将所有筹码都押在一段关系上；让对方感到失去你是有代价的。

\subsection{平衡的艺术：骨头与鞭子的动态协调}

\subsubsection{根据阶段调整比例}

关系不同阶段，骨头与鞭子的比例应有所不同。早期可能需要更多给予以建立信任，中期需要更多约束以巩固边界，后期需要更多平衡以维持稳定。

\textbf{实践建议}：早期80\%给予，20\%约束；中期60\%给予，40\%约束；成熟期50\%给予，50\%约束。

\subsubsection{根据对象调整策略}

不同的人对给予和约束的反应不同。有的人需要更多安全感，有的人需要更多挑战。根据对方的性格和需求调整策略。

\textbf{实践建议}：安全型人格平衡给予和约束；焦虑型人格更多给予和确认；回避型人格更多空间和自主。

\subsubsection{根据情境灵活应变}

关系不是静态的，需要根据情境灵活调整。在对方经历困难时可能需要更多支持，在对方自满懈怠时可能需要更多压力。

\textbf{实践建议}：敏感捕捉对方的状态变化；不固守某一策略；与对方保持沟通。

\section{可能的批评与回应}

\subsection{批评一：这一类比是否将亲密关系过度政治化？}

\textbf{批评}：亲密关系是情感和道德的领域，不应该用权力逻辑来分析。将马歇尔计划的权力技术应用于恋爱，是将工具理性入侵到了本应纯粹的情感领域。

\textbf{回应}：这一批评有一定道理，但忽视了权力的普遍性。权力不是只存在于政治领域，而是存在于一切社会关系之中。承认亲密关系中的权力维度，不是要将爱情变成算计，而是要正视现实，以更成熟的方式经营关系。

正如福柯所言，权力不是邪恶的，它是一切社会关系的构成要素。关键不是否认权力的存在，而是让权力以健康、建设性的方式运作。

\subsection{批评二：这一模型是否忽略了爱情的独特性？}

\textbf{批评}：爱情包含激情、牺牲、无条件付出等要素，这些都是马歇尔计划模型无法涵盖的。用政治逻辑分析爱情，错过了爱情的本质。

\textbf{回应}：爱情确实包含这些要素，但长期稳定的关系不止有爱情。激情会消退，无条件付出可能耗尽。可持续的关系需要某种结构性的支撑，而这正是权力动力学可以解释的层面。

本文不否认爱情的独特性，只是指出：仅有爱情不足以维持长期关系。权力的健康运作是关系持久的结构性条件之一。

\subsection{批评三：这一模型是否鼓励操纵？}

\textbf{批评}：教导人们如何运用"骨头与鞭子"，可能鼓励在亲密关系中进行操纵和控制。这在道德上是可疑的。

\textbf{回应}：这是一个严肃的担忧。确实，本文描述的权力技术可以被滥用。但本文的目的是分析权力的运作逻辑，而非鼓励操纵。

更重要的是，本文强调的是"建设性的权力"——权力服务于双方的成长和关系的健康，而非单方面的利益。马歇尔计划之所以成功，正是因为它最终造福了西德和日本，而非仅仅是控制了它们。

同样，在亲密关系中，健康的权力运作应当让对方变得更好，而非让对方变得依赖和脆弱。这是"马歇尔计划式"关系与操纵性关系的根本区别。

\section{结论：权力的温柔与残酷}

\subsection{主要发现}

本文通过分析马歇尔计划这一历史案例，提炼出一套"骨头与鞭子"的权力运作逻辑，并将其应用于亲密关系分析。主要发现包括：

\begin{enumerate}
    \item \textbf{结构性依赖是持久关系的基础}。无论是国际关系还是亲密关系，让对方在某些关键维度上依赖你，是关系持久的前提。
    
    \item \textbf{给予与约束必须平衡}。单纯给予会创造依赖但可能失去尊重，单纯约束会制造恐惧但破坏认同。持久的权力关系需要两者的动态平衡。
    
    \item \textbf{合法性来自正向体验}。权力需要合法性，合法性来自让对方体验到合作的好处。经济增长为马歇尔计划提供了合法性，幸福成长为亲密关系提供了合法性。
    
    \item \textbf{边界控制优于全面控制}。过度控制会窒息活力，明智的控制是设定边界并在边界内给予自由。
    
    \item \textbf{长期主义是关键}。权力关系的建设需要时间，短期机会主义会破坏长期基础。
\end{enumerate}

\subsection{理论贡献}

本文的理论贡献在于：建立了国际政治与亲密关系之间的跨学科分析框架；提出了"战略性依赖"和"建设性张力"两个核心概念；为理解亲密关系中的健康权力运作提供了新的理论视角。

\subsection{实践启示}

对于亲密关系实践者，本文提供了以下启示：识别并满足对方的核心需求，创造不可替代的价值；设定健康边界，在边界内给予充分自由；以促进成长为导向施加压力，避免为控制而控制；持续创造正向体验，为关系生产合法性；采取长期主义视野，投资于关系的可持续性。

\subsection{最后的话}

权力是温柔的，因为它创造依赖、提供保护、促进成长；权力也是残酷的，因为它设定边界、施加压力、维持不对称。马歇尔计划展示了权力的两面性：它既是馈赠，也是控制；既是解放，也是约束。

亲密关系中的权力同样如此。健康的权力不是否认或逃避，而是正视并以建设性的方式运用。让对方依赖你但不压垮对方，约束对方但不窒息对方，引导对方但不替代对方——这是亲密关系中的"马歇尔计划"，也是爱的政治艺术。

在某种意义上，所有的长期关系都是一场微型政治。理解这一点，不是要让爱情变得冷酷，而是要让爱情变得清醒。只有在清醒中，我们才能真正地爱——不是盲目的、自我牺牲的爱，而是成熟的、相互成就的爱。

这或许就是"骨头与鞭子"的终极启示：权力的最高境界，是让被统治者心甘情愿地追随，让对方在依附中获得成长，在约束中体验自由。

这才是真正的驯化艺术——无论是国家还是心灵。

\newpage
\begin{thebibliography}{99}

\bibitem{marshall1947}
Marshall, G. C. (1947). \textit{Commencement Address at Harvard University}. Harvard University Archives.

\bibitem{maier1981}
Maier, C. S. (1981). \textit{The Two Postwar Eras and the Conditions for Stability in Twentieth-Century Western Europe}. American Historical Review, 86(2), 327-352.

\bibitem{hogan1987}
Hogan, M. J. (1987). \textit{The Marshall Plan: America, Britain, and the Reconstruction of Western Europe, 1947-1952}. Cambridge University Press.

\bibitem{dower1999}
Dower, J. W. (1999). \textit{Embracing Defeat: Japan in the Wake of World War II}. W. W. Norton \& Company.

\bibitem{gaddis2005}
Gaddis, J. L. (2005). \textit{The Cold War: A New History}. Penguin Press.

\bibitem{bowlby1969}
Bowlby, J. (1969). \textit{Attachment and Loss: Vol. 1. Attachment}. Basic Books.

\bibitem{sternberg1986}
Sternberg, R. J. (1986). \textit{A Triangular Theory of Love}. Psychological Review, 93(2), 119-135.

\bibitem{foucault1978}
Foucault, M. (1978). \textit{The History of Sexuality, Volume 1: An Introduction}. Pantheon Books.

\bibitem{gottman1994}
Gottman, J. (1994). \textit{What Predicts Divorce? The Relationship Between Marital Processes and Marital Outcomes}. Lawrence Erlbaum Associates.

\bibitem{baumeister1995}
Baumeister, R. F., \& Leary, M. R. (1995). \textit{The Need to Belong: Desire for Interpersonal Attachments as a Fundamental Human Motivation}. Psychological Bulletin, 117(3), 497-529.

\end{thebibliography}

\end{document}
